\documentclass{ltjsbook}
\usepackage{docmute} 
\usepackage{../../myStyle} 

\begin{document}
\chapter{プログラミングの基礎}
\label{x16_Programming}
さて，ここから始まる統計の話は，統計学の中でも最近使われるようになってきたベイズ統計を駆使し，色々な心理学的現象を分析していこうというものになります。

このアプローチはデータに合わせて分析モデルを考えていくことになるので，ボタンをポチポチ押すと答えが出るような定型パターンの分析ではありません。個別のケースに合わせて分析モデルを考えていくことになります。その時に必要になってくるのがプログラミング技術ですので，まずはウォーミングアップということでプログラミングの話を始めていきたいと思います。中にはうんちくのようなところもありますので，興味があるところだけ飛ばして読んでいただいて構いません。最後にあげた課題ができるのであれば，知識は後からでもいいと思います。

\section{プログラミングの基礎}
プログラミングとは，コンピュータに計算をさせるその仕様書を作ること，だと思ってください。お料理のレシピのように，計算レシピを書くわけですね。ただ相手は計算機ですので，言葉の端々まで正確に伝えないと理解してくれません。よく「思った通りに動いてくれない！」と不満を訴える初心者がいますが，それは当然で，プログラムは\textbf{思った通りに動くのではなく，書いた通りに動く}のです。思った通りに動かないのであれば，それは仕様書・レシピの方に間違いがあります。

今の第一原則に加えて，プログラミングを進める上での注意点をあげておきます。簡単なことだと思うかもしれませんが，基本的なルールをしっかり守ることが上達への近道です。
\begin{enumerate}
	\item プログラムは思った通りには動かない。書いた通りに動く。大文字，小文字，スペルの違いに注意する。
	\item 書き間違えないための工夫は美しさ。綺麗に書くことが大事。
	\item 一言一句すべてに意味がある。ただの写経ではなく，意味を考えながら書く。
	\item 「遊び心」をもって！ここを変えたらどうなるか，を少しずつ「やってみる」が大事
	\item 変更は少しずつ。一気に変えるとどこが変わったかわからなくなるから。
\end{enumerate}

\subsection{綺麗に書こう}
綺麗に書く，というのはコードの可読性を上げるために必要です。綺麗な書き方として，\citet{松浦201610}は次の5点をあげています。

\begin{itemize}
	\item インデントは必ずする
	\item データを表す変数の先頭の文字は大文字。パラメータを表す変数の先頭の文字は小文字。
	\item 各ブロックの間は１行あける
	\item camelCaseではなく\verb|snake_case|で。
	\item 「~」や「＝」の前後は1スペースあける
\end{itemize}

1つめのインデントとは，字下げのことです。行の頭を少し凹ませることで，違う人まとまりであることを明示するのです。たとえば次の2つのコードは同じ働きをしますが，\ref{code::16_02}の方が見やすいと思いませんか？
\begin{lstlisting}[language=c++,label={code::16_01},caption={インデントのないコード}]
data {
int<lower=0> N;
vector[N] y;}
parameters {
real mu;
real<lower=0> sigma;}
model {
for(i in 1:N)
y[i]~normal(mu, sigma);}
}
\end{lstlisting}


\begin{lstlisting}[language=c++,label={code::16_02},caption={インデントのあるコード}]
data {
  int<lower=0> N;
  vector[N] y;
}

parameters {
  real mu;
  real<lower=0> sigma;
}

model {
    for( i in 1:N ){
        y[i] ~ normal(mu, sigma);
    }
}
\end{lstlisting}
インデントのある\ref{code::16_02}は，\texttt{data}というブロック(\{ と\} で囲まれている領域)の中に，2つの行が入っている，ということが明確にわかります。また\texttt{model}というブロックの中には，\texttt{for}から始まる分がありますが，これも複数行に渡るブロックを構成する文なので，その中身はインデント(字下げ)されています。このように，インデントすることでどこの列がどこまでブロックを組んでいるのかがわかりやすくなるのです。ちなみにこのインデントを作るのはTABキーを使います。TABキーは普段，どう使うのかわからないものだと思われがちですが，このインデントをしてくれるためのものです\footnote{ここには流派があって，TABでインデントする派とスペースキーを4回または8回押してインデントする派がいます。どちらでも働きは同じなのですが，個人的には数回の字下げを1回のキーで行ってくれるTABのほうが良いと思っています。}。コードエディタによっては，カッコで括ったときに自動的に閉じるカッコを用意し，また改行ごとに字下げ位置を合わせてくれるものがあります。これらの機能を使って是非わかりやすいコードを書いてください。

\citep{松浦201610}の指摘の2,3番目についてはお好みで，4番目もそれほど強い・広く浸透した決まりではありません\footnote{ちなみにキャメルケースCamel Case，スネークケースSnake Caseとは変数名のつけ方のルールです。Rなどプログラミング言語ではあらゆるものに名前をつけて管理しますが，名前のつけ方は任意です。命名はわかりやすい方がいいですが，長いものになると区切りを入れたくなるかもしれません。さてCamelとはラクダの意味で，ラクダのコブのように区切りのところだけ大文字にする，という記法です。Snakeは蛇のことで，区切りのところにアンダースコア\verb|_|を使うというものです。たとえば野球チーム，という変数名をつけたい時に，キャメルケースのやり方だと\texttt{BaseballTeam}のように書きますし，スネークケースのやり方ですと\verb|Baseball_team|のような書き方になります。ちなみにRは日本語での命名も許しますが，そのほかの言語では一般的ではありませんし，何より全角と半角を切り替える時のミスやエラーが多くなるので，半角英数字での命名をすべきという点はどちらでも共通です。}。ただし，5番目のスペースの前後に少し空白を取るのは，見やすさのためにも是非実行してもらいたいところです。

これらの書き方は，\keytermL{可読性}{かどくせい}を上げるためのものです。プログラムは仕様書・レシピですから，あとで読み直した時やほかの人が読むときも，意味がわかるようにしておいた方が良いのです。自分だけわかれば良い，と思うかもしれませんが，その自分ですら何があるのかわからなくなってしまう可能性があります。そのような読みにくさはすぐにバグにつながりますから，綺麗に書くことを常々心がけておいて欲しいのです。

\subsection{意味を考えて書こう}
プログラミング言語は，ほとんど英単語のようなものです。プログラミング上の都合から，短い言葉に省略されていたり，アンダースコアやピリオドでつながって書いてあったりしますが，比較的わかりやすい言葉であることに違いはありません。

プログラムを書く上で，最初は写経と呼ばれる，テキストや指示に従って書き写すことが基本です。もちろんネット上にあるサンプルコードなどをコピー＆ペーストしても良いのですが，自分用の分析コードを書くためには手を入れる必要があったりします。ですから，最初はなるべく自分でキーボードを叩いて，コピー＆ペーストではなく入力する経験を積んで欲しいです。もちろん間違えてしまうことが多いでしょうが，転ばずに歩けるようになった人が一人もいないように\footnote{お釈迦様は除く。}，ミスを犯すことでどういうミスだったかを考え，自覚して，ミスが減っていきます。失敗する経験も時には必要なのです。

ということで，時折自分なりにコードを書いてもらうのですが，その時に「ただの記号列」と思わないでください。既に書いたように，英語あるいはそれを省略した表現になっているので，見慣れないものかもしれませんが必ず意味があります。
プログラムを始める最初の頃は，ミススペルが多く，あちこちでエラーが出て気持ちが萎えることもあるかもしれません。エラーが出る時は目を凝らして，どこに問題があるかを考えて修正することになります\footnote{ちなみにプログラミング歴30年以上の私でも，簡単な英単語，たとえばlibraryのスペルを間違えるなんてことはしょっちゅうです。}。ここで難しいのが，$x$や$s$などの大文字と小文字の区別がつきにくいもの，$1$と$l$のように違いがわかりにくいものがあるということです。
私は職業柄，多くの人に教え，多くのバグを見つけてきましたが，その中でも特別見つけるのに苦労したのが，\texttt{norma1}というミススペルでした。みなさん，これのどこがミススペルかわかりますよね？でもこれ，書く時に「正規分布のことだな」と思っていれば，そもそも入力時にそんな入れ方はしないと思うのです。ただの字だ，と心を無にしてしまうと，後々見つけにくいエラーを作ることにもなりますから，この文字・この名前・この関数はどういう意味だろう，と少し考えながら取り組んでみてください。

\subsection{遊び心が大事}
プログラミングを進める上では，遊び心が大事です。うまく動くコードがかけたら，もうおかしなことになりたくない，と手をつけないままにしてしまう人がいます。

ところで，今まで教えてきた経験から言って，プログラミングが上達する人は，うまくコードが書けたらすぐに「ここをこう変えたらどうなるんだろう」と試してみる人だと思います。最初はプログラムの意味が全くわからないので，どこをどういじっていいのか見当もつかないかもしれません。しかしたとえば\texttt{blue}という文字が出てきたら，これは色の青のことじゃないか，と思いますよね。ではその\texttt{blue}を\texttt{red}に変えたらどうなるんだろう？やってみよう！となるような，そういう遊び心が大事なのです。

もちろん変えてしまったことでエラーになって，動かなくしてしまうことがあるかもしれません。その場合は，元に戻して元通り動くかどうか，さらに確認すれば良いのです。数字や文字を少しいじっただけで，PCが爆発してしまうようなことはありませんから，ちょっと遊び心を出して，どうなるのかな？と思う気持ちを大事にしてください。

\subsection{変える時は少しずつ}
プログラムを色々いじって遊ぶことが成長につながる，という話をしました。ただし遊び方には気をつけましょう。まずうまくいくコードができれば，それは保存しておいて，また同じ内容のものを別名で保存し，「遊んでいい方」「壊れてもいい方」を作って，そちらで色々変えて遊ぶといいでしょう。最悪なことがあっても，うまくいくバージョンは常に残っているわけですから。

そして色々試す時のコツは，一箇所ずつ変更していくことです。たとえば\texttt{blue}を\texttt{red}に変え，\texttt{line}を\texttt{box}に変え・・・と複数箇所を同時に変更して，うまくいかなかった場合，どちらが原因担っているのか把握できませんよね。これは心理学実験で言うところの交絡というやつです。操作が交絡してしまって，原因が特定できなくなるのです。

また，実行するときも1行ずつやりましょう。特にRは一問一答型，つまり1つの命令について1つの反応を随時返してきます。複数行をまとめて実行することもできますが，まとめて実行している途中でエラーが出ていて，その後の計算が全て空回りしていると言うことも少なくありません。エラーがどこで生じているのか，しっかり特定することが重要です。そのためにも焦らず，1つずつ確実にできることを積み重ねていくという姿勢が重要です。

統計分析も複雑で細部まで設定し尽くしたモデルを作り上げていくことができますが，いきなり全体が完成するのではなく，小さなピースの積み重ねなのです\footnote{Rやプログラミングで統計を学ぶ意義はここにあります。GUIで操作できる統計パッケージも少なくありませんが，それらは画面が出てきた時にデフォルトで設定されている値があるのがほとんどで，自分が何をやっているのか細部まで気づかないまま実行できてしまうのです。そうすると当然，エラーが出たり，うまくいっても誤用している，と言うことにもなりかねません。自分で理解できていない技術に振り回されないようにするために，しっかりとわかるピースを積み重ねていく必要があるのです。}。


\section{プログラミング言語}

プログラミングはコンピュータを動かすための仕様書を書くことです。ここで「コンピュータ」というのは計算機という意味だと思ってください。もちろんコンピュータは数字の計算だけでなく，音楽を奏でたりゲームを楽しませてくれたり，と色々なことができるのですが，その背後にあるのはとにかく0/1の数値演算です。0か1かという数字がどうして映像や音声，通信になるのか，と疑ってしまうほど異なっているようですが，それでも全て数字の計算から構成されています。

コンピュータができ始めたごく初期の頃は，これを動かすのに0と1からなる数字の羅列で仕様書を用意していました。流石にそれではわかりにくく表現力にも乏しいので，次にできたのがマシン語とよばれる16進数による表現でした。この段階でも，知らない人にとっては暗号か，意味のある連なりに思えない文字列にすぎません。画面に線を引いて欲しい時に\texttt{line}という命令で伝えられるようになって，やっと普通の人間にも意味がわかるレベルになってきます。このように，人間がわかるレベルでコンピュータに命令が伝えられるような言語のことを\keytermL{高級言語}{こうきゅうげんご}と言います。高級言語は人間よりで，直接コンピュータが理解できませんので，この高級言語を機械の言葉に翻訳するためのアプリケーションを介在させます。それがいわゆるプログラミング言語と呼ばれるものです。

プログラミング言語は，古くはBASIC，PASCAL，FORTANなどと呼ばれる書式のものがあり，その後出てきたC言語，その改良版であるC++言語\footnote{この\verb|++|という書き方は，C言語特有の「変数に1加える」という表記法であり，C言語の改良版という意味でC++と命名されています。読み方はシープラスプラスですが，シィプラプラの愛称でも知られています。}などが有名です。他にもJAVAやObject C，Pythonなどが有名ですね。Rも言語の一種で，統計解析に特化したものです。また，この授業で扱うStanという確率型プログラミング言語は，確率モデルの分析に特化したものです。みなさんがデータサイエンス業界に進むのであれば，PythonやRを使えるようになっておくといいでしょう。これらの言語の特徴の1つは，パッケージを追加することで機能が拡張できることにあります。特に機械学習系(AIなど)のパッケージが豊富なのがPythonです。またこれらの言語は基本的にフリーソフトウェアであり，導入にお金がかからないところもポイントですね。タダで始められるおもちゃみたいなものです！\footnote{余談ですが，Scratchや任天堂Switchの「ナビつき! つくってわかる はじめてゲームプログラミング」などにみられる，GUIベースのプログラミング言語もあります。これはプログラミングの要素がアイコンやキャラクターになっていて，それらを組み合わせることで計算させるという方法をとっています。その操作のしやすさ(ミススペルがない！)から，小学生などにも導入できる素晴らしい取り組みです。同様の発想は，LOGO言語などにもみられました。}\footnote{PerlやJava Scriptなど，Webブラウザ上で動くプログラミング言語もありますが，これはブラウザ上での操作に限定されており，数値計算にはあまり向かないのでここには取り上げていません。先程あげたJAVAとJAVA Scriptは別物なので注意してください。}\footnote{GUIアプリケーションを作ることに特化した言語もあります。数値計算と違って，マウスがどこでクリックされたか，と言った「動き」を契機としてプログラムが走る，イベントドリブン型の言語で，Microsoft社のVsial Basicなどが有名です。これはMicrosoft Office製品の中に埋め込むこともでき(VBA)，これを使うとExcelでも高度な統計解析ができたりします。Excelを使った統計ソフトウェアとして代表的なものにHADがあります\citep{HAD}。他にもDelphiなどの言語がありました。}

これらの高級言語は，完全な日本語・英語ではありませんが，略語や比較的意味のわかる形で計算仕様書を書かせてくれます。各言語は我々の書く仕様書をコンピュータがわかる0/1の形に翻訳して伝え，複雑な計算を実行してくれるのです。この翻訳の仕方として，Rのようにひとつひとつの命令文を逐次的に実行する一問一答型の方式と，仕様書全体をまとめて翻訳し，結果を一気に返す方式の2つのパターンがあります。前者を\keyterm{インタプリタ}{いんたぷりた}{interpreter}型，後者を\keyterm{コンパイラ}{こんぱいら}{compiler}型と言います。インタプリタ型は命令に対する回答がすぐに返ってくるので，どこでエラーが出たのかがわかりやすいという利点があります。しかし毎回その場で同時通訳するので，全体的なスピードが遅いという欠点があります。コンパイラ型は仕様書をまとめて翻訳するので，文書全体を翻訳し終わって実行させますから計算スピードが早いのですが，翻訳にちょっと時間を要すること，エラーが発生したらコンパイルをやり直す必要があることなどが欠点です。ちなみにコンパイルされたものが実行ファイルであり，おおくのアプリケーションはこの翻訳済みの実行ファイルの形で提供されます。実行ファイルは機械語で書かれていますから，人間が読んでもわかりませんが，機械にあらゆる動きをさせることができるものです。機械に致命的な故障を生じさせるバグがあったり，悪意ある操作が含まれていても実行ファイルを一瞥しただけではわかりませんから，ウイルス対策ソフトなどはコンパイラを危険なソフトウェアと判断することもあります。

この授業で使うR言語はインタプリタ型，Stanはコンパイラ型です。この両者の違いも意識しておいた方が良いかもしれません。

\section{高級言語の基本的な働き}
プログラミング言語は様々ありますが，それらのやっていることは基本的に計算であり，その表現方法が違うだけです。どの言語にも共通する，プログラミング言語の特徴的な働きは「代入」と「反復」と「条件分岐」です。以下，この3つの働きについてRを使ってみていきます。
\subsection{代入}
料理でよくあるシーンとして，「野菜を切ってザルにあげておく」など，一旦作業の途中経過を横に退けて別の作業をする，というのがあります。計算プロセスでも，「計算結果を一旦横に置いておく」という操作をしたいことがあります。これが\textbf{代入}で，コンピュータとしてはあるメモリ番地に値を書き込んでおく，という操作をすることになります。横に置いておく時に名前をつけることができ，これがRのなかでは\keyterm{オブジェクト}{おぶじぇくと}{Object}と呼ばれるものになります。すでにみなさんもやったことがあると思いますが，計算結果をオブジェクトに代入しておくコードの例をみてみましょう(\texttt{code::\ref{code::16_12}})。

\begin{lstlisting}[language=c++,label={code::16_12},caption={代入操作}]
a <- 1
b <- 2
a + b
a <- 3
a + b
\end{lstlisting}

\paragraph{コード解説}
\begin{description}
	\item[1行目] オブジェクト$a$に$1$を代入
	\item[2行目] オブジェクト$b$に$2$を代入
	\item[3行目] オブジェクト$a$と$b$に$+$という操作。すなわち$1+2$の計算を指示していることと同じ。
	\item[4行目] オブジェクト$a$に$3$を代入。もともと入っていた情報は上書きされる。
	\item[5行目] オブジェクト$a$と$b$に$+$という操作。すなわち$3+2$の計算を指示していることと同じ。
\end{description}

非常に単純な例で，今更何を，と思った人もいるかもしれませんが，これが基本中の基本なので改めて示しました。ポイントは3行目の「オブジェクト名で値を指定していること」と，4行目の「値の上書き」にあります。オブジェクト名で計算を表現することで，操作の一般化ができます。すなわち，3行目と5行目が同じ式であるのに結果が違うように，「どんなに値が違っても同じ動作をする」ようにできるわけです。コンピュータがすごいのは，この手の代入がどれほど複雑に，どれほど繰り返されても決して計算ミスをせず，指示通りに動くことができるというところです。人間は「ベクトルの2番目の要素と3番目の要素を足す」という操作を数回やるだけでも計算ミスをすることがあり得ますが，コンピュータは同じ計算を100回やっても1000回やってもミスなく動作します！

\subsection{反復}
エビフライをたくさん作ることを考えてみましょう。「えびの殻を剥いて，衣をつけてあげる」という動作を何回も繰り返すことになると思います。このような反復捜査もコンピュータは得意とするところです。なんせ，疲れを知りませんので。
反復操作を指示するコードは\texttt{for}文という書式を使います。反復計算コードの例は次のようなものです(\texttt{code::\ref{code::16_22}})。

\begin{lstlisting}[language=c++,label={code::16_22},caption={反復操作1}]
a <- 0
for (i in 1:4) {
  a <- a + 1
  print(a)
}
\end{lstlisting}

\paragraph{コード解説}
\begin{description}
	\item[1行目] オブジェクト$a$に$0$を代入(初期化)
	\item[2行目] \texttt{for}文開始。中括弧で括られている5行目までの中身を反復計算する。
	\item[3行目] オブジェクト$a$に，もとの$a$に$1$を加えた値を代入(上書き)している。
	\item[4行目] オブジェクト$a$を出力させている。
	\item[5行目] \texttt{for}文ここまで。
\end{description}

ここでは3行目の「1を加える」，4行目の「表示する」という操作を繰り返していることになります。
ポイントは2行目の書き方ですね。\texttt{for}の後ろの小カッコ()の中身が，繰り返しに使う要素の指定です。後ろの$1:4$はR言語特有の書き方で，「1から4まで」すなわち\verb|1,2,3,4|を意味しています。変数$i$がこの順番で変わるよ，ということを意味しているので，$i$は1回目，2回目，3回目，4回目，と進んでいくことになります。\texttt{code::\ref{code::16_03}}のような書き方をすると，その振る舞いがよくわかるかと思います。

\begin{lstlisting}[language=c++,label={code::16_03},caption={反復操作2}]
a <- 0
for (i in c(1, 3, 5, 15, 12)) {
  print(paste(a, "に", i, "を加えます"))
  a <- a + i
  print(a)
}
\end{lstlisting}

ここでのポイントは，反復用インデックス$i$が$1,3,5,15,12$の順に変わっていくということと，変わっていくインデックス$i$の値そのものも計算に使えるということです。2点目については注意が必要で，反復用インデックスが$1,2,3,4$と変わるような例の場合，途中で$i <- 1$などとしてしまうと，いつまでもインデックスが前に進まず永遠に計算が終わらないことになります。そのようなミスが生じないように，注意してくださいね。

また，この\verb|for|文は入れ子にして使うことができます。例えば$i$が1から5まで変わり，$j$が1から3まで変わりながら計算をする，ということを考えて，次のような表現が可能です。
\begin{lstlisting}[language=c++,label={code::16_32},caption={反復操作1}]
A <- matrix(1:15, ncol = 3, nrow = 5)
A
for (i in 1:5) {
  for (j in 1:3) {
    A[i, j] <- A[i, j] + i * j
  }
}
A
\end{lstlisting}

\paragraph{コード解説}
\begin{description}
	\item[1--2行目] オブジェクト$A$は3行5列の行列で，1から15までの数字が順に入っている。
	\item[3行目] \texttt{for}文開始。$i$が1から5まで変わる。
	\item[4行目] \texttt{for}文その2開始。$j$が1から3まで変わる。
	\item[5行目] 行列$A$の$i$行$j$列目の要素に対し，$i\times j$の計算結果を加えたものを代入(上書き)させる。
	\item[6行目] \texttt{for}文その2がここまで。
	\item[7行目] \texttt{for}文その1がここまで。
	\item[8行目] 最終計算結果の表示。
\end{description}

行列の行，列それぞれについて順番に，各要素を指定しながら代入していくという計算です。挙動を確認しておきましょう。

\subsection{条件分岐}
\label{if_statement}
「サラダ油が200円より安かったら買ってくる」というようなお使い指示，ありますよね。これは「200円以上の時は買わない」という意味でもあります。この200円かどうか，という条件に対して，その後の動きが「買う」「買わない」に分岐するので，条件分岐と呼ばれます。これをプログラムで書くと次のようになります。
\begin{lstlisting}[language=c++,label={code::16_04},caption={条件分岐}]
egg <- 250
if (egg < 200) {
  print("卵を買います")
} else {
  print("卵を買いません")
}
\end{lstlisting}

\begin{description}
	\item[1行目] オブジェクト$egg$に$250$を代入。
	\item[2行目] \texttt{if}文開始。()内が条件で，条件が該当すれば次の中括弧で括られている領域までを実行する。
	\item[3行目] 「卵を買います」と画面表示させる。
	\item[4行目] 条件が該当した時の実行内容を閉じつつ，該当しない場合の実行内容を書く領域を展開する。
	\item[5行目] 「卵を買いません」と画面表示させる。
	\item[6行目] 条件が該当しない場合の実行内容を閉じる。
\end{description}

一行目の\verb|egg|に様々な数字を入れて検証してみてください。思った通りの振る舞いができているでしょうか。
この\texttt{if}文は小括弧の中身が条件節ということになります。条件が成立していることを\textbf{真}または\textbf{TRUE}と表現し，成立しないことを\textbf{偽}または\textbf{FALSE}と表現します\footnote{\texttt{TRUE,FALSE}が大文字なのは重要で，真偽を表すRの特別な用語です。こうした用語のことを予約語と言います。}。今回は\verb|egg|変数が200未満であることを条件としています。もし200円も含めたいのであれば(200円以下にしたいのであれば)，\verb|egg <=200 |と書く必要があります。

条件分岐をする場合は，条件が思った通りに設定できているか，分岐した後のルートが間違いなく書かれているかなどに注意が必要です。というのも，我々が日常言語で使っている「もし〜ならXXXする」といった表現は曖昧なことがあり，書いた通りにしか解釈しないコンピュータは，思った通りに動いてくれないということがあるからです。Twitter上で有名になったジョークに，次のようなものがあります\footnote{\url{https://twitter.com/beamtetrode350b/status/1406773935069229057}}。

\begin{quote}
	ある妻がプログラマの夫に「買い物にいって牛乳を1つ買ってきてちょうだい。卵があったら6つお願い」と言った。
	夫はしばらくして、牛乳を6パック買ってきた。
	妻は聞いた「なんで牛乳を6パックも買ってきたのよ！」
	夫いわく「だって、卵があったから……」
\end{quote}

これはプログラミング的思考の例として示されており，気の利かない夫とされていますが，言外の意味を含めすぎた妻の条件分岐指示がまずかったとも言えます。より適切な指示の例は，次のようなものです\footnote{\url{https://twitter.com/takl/status/1127065591380971521}}。

\begin{quote}
	プログラマを夫にもち長年苦労してきた妻「1パック8個以上入った190円以上210円以下の鶏卵のパックがあったら買ってきて。それがなかったら1パック8個以上入った189円を超えない最も高い鶏卵のパックを買ってきて。売り場について15分以内に見つけられなかったときは私に電話してから指示に従って」
\end{quote}

これでもまだ指示が厳密でない！条件漏れの可能性がある！というコメントがありますが，「書いた通りにしか動かない」ことの重要さがお分かりいただけるかと思います。


\section{まとめ}
ここで説明したのは，レゴのピースのようなものです。ピース1つ1つは小さくて，それだけではなんの形も作り上げることができませんが，細かなピースでも組み合わせ次第で大きなオブジェを作ることができます。
プログラミングにはたった1つの正解というのはなく，どういう形であれいと通りに動くものができればそれで構いません。書き方は人それぞれでよく，結果が伴うかどうかがポイントです。
どのピースをどのように組み上げてもいいので，思い通りのものが作れるようにトレーニングをしておきましょう。

\section{課題}
\paragraph{FizzBuzz課題}
1から15までの数列に対して，次の条件にあった出力をするプログラムを書きなさい。なお，割り算の余りを計算するRの関数は\verb|%%|，文字を出力する関数は\verb|print|です。
\begin{itemize}
	\item その数が3で割り切れる場合には、その数の代わりに「Fizz」を出力する
	\item その数が5で割り切れる場合には、その数の代わりに「Buzz」を出力する
	\item その数が3でも5でも割り切れる場合には、その数の代わりに「FizzBuzz」を出力する
\end{itemize}

\paragraph{行列のかけ算}
行列$\bm{A} = \begin{pmatrix}1&2&3\\4&2&0\end{pmatrix}$と$\bm{B}=\begin{pmatrix}1&5\\2&3\\3&8\end{pmatrix}$を用意し，$\bm{AB}$の掛け算をしたいとします。
Rの行列の積を求める記号ではなく，\verb|for|文をつかって，この計算をするプログラムを書きなさい。

\end{document}
